\section{Introduction}
\label{sec:intro}

Post-training quantization (PTQ) converts a trained network to low-precision weights without a full training stage~\cite{gholami2021survey}.
Research on PTQ mostly improves the quantizer, through better calibration~\cite{yuan2022ptq4vit}, mixed precision~\cite{kim2026lampq} or synthetic calibration data~\cite{zhong2022fdda,kim2025synq}, and evaluates it on a few standard benchmarks.
However, the same quantizer can cost one dataset far more accuracy than another.
Degradation varies across medical imaging datasets and backbones~\cite{ozdemir2026bridging}, across classes and demographic groups~\cite{liu2025unfairness}, and with the calibration distribution~\cite{yuan2023benchmarking}.
At the sample level, prediction changes after quantization have been linked to small output margins~\cite{hu2023towards}.
These findings make quantization robustness a property of the data and the trained model, not only of the quantizer.

\begin{figure}[t]
  \centering
  \includegraphics[width=\linewidth]{figures/retention_curves.pdf}
  \caption{Accuracy retention against weight bit width at group size 128, averaged over three backbones and three checkpoint regimes. Thin lines are datasets, thick lines group means. VegFru is excluded because its grid is incomplete.}
  \label{fig:teaser}
\end{figure}

Fine-grained image recognition (FGIR) makes this dependence visible, because its datasets range from birds and cars to soybean and cotton cultivars that experts struggle to tell apart~\cite{yu2021benchmark}.
\cref{fig:teaser} shows that weight-only quantization treats them unevenly.
At 3-bit weights, each of the five cultivar datasets keeps at most \V{maxSc_w3g128} of its accuracy, while each other dataset keeps at least \V{minOther_w3g128}.
However, the cultivar datasets also have much lower baseline accuracy, so final accuracy alone cannot say which dataset property matters.
We therefore ask where in the bit-width and group-size grid the dataset effect appears, whether the ranking of datasets holds across models, and whether baseline decision margins account for the difficulty beyond baseline accuracy.

To answer these questions, we compare every quantized model with its own unquantized checkpoint over a full grid of weight-only quantization conditions~\cite{badri2023hqq}.
Our contributions are as follows:
\begin{enumerate}
  \item We measure accuracy retention for \V{nDatasets} fine-grained datasets, three backbones and three checkpoint regimes over \V{nConditions} conditions, and locate the dataset effect in the transition band between collapse and full retention.
  \item We show that dataset difficulty rankings are stable across quantization conditions and only partly stable across backbones and regimes.
  \item We test prespecified margin diagnostics and show that baseline accuracy explains much, but not all, of the difficulty.
\end{enumerate}
